% !TeX root = ../复分析笔记.tex
% 由主文档通过 \input 载入；请编译主文档。
\setcounter{chapter}{13} % 保留原书章号。
\chapter{解析延拓与黎曼曲面初步}
\label{chap:ca-continuation}

泰勒级数只在一个圆盘中表示函数，但函数本身可能越过这个圆盘继续存在。例如 $\sum_{n\ge0}z^n$ 的收敛半径为一，而 $1/(1-z)$ 在 $\mathbb C\setminus\{1\}$ 上全纯。解析延拓研究怎样由局部信息确定更大范围的函数；路径绕行又说明，局部唯一性未必自动产生全局单值性。本章把这两件事分别处理，并用具体模型引出黎曼曲面。

\section{函数芽与相容的局部函数}

\begin{definition}[函数芽与直接解析延拓]
设 $f$、$g$ 分别在点 $a$ 的某些开邻域上全纯。若它们在 $a$ 的一个更小邻域上相同，就称二者在 $a$ 定义同一个\emph{全纯函数芽（holomorphic germ）}，记为 $[f]_a=[g]_a$。若 $f$ 在区域 $U$ 上全纯，$g$ 在更大区域 $V\supset U$ 上全纯且 $g|_U=f$，则 $g$ 是 $f$ 的解析延拓。更一般地，两个相交区域上的函数可以在某个非空交叠开集上相同，并借此向新的区域延续。
\end{definition}

函数芽保留一点附近的全部函数信息，而不固定一个过大的定义域。相等关系具有自反、对称、传递性，因此确实给出等价类。恒等定理说明：在同一个连通延拓区域上，若延拓存在，则至多一个。但是，若两个定义域的交集不连通，在其中一个分支上相同不强迫在其他分支上也相同；粘合单值函数时必须在全部交集上相容。

\begin{example}[换中心扩展几何级数]
原点处 $f(z)=\sum_{n\ge0}z^n$ 在 $|z|<1$ 内等于 $1/(1-z)$。在 $a=-1/2$ 附近，
\[
 \frac1{1-z}=\frac{2}{3}\sum_{n=0}^{\infty}
 \left(\frac{2(z+1/2)}3\right)^n,
 \qquad |z+1/2|<3/2.
\]
两个圆盘的交集非空，且两式同为同一个有理函数，所以第二个级数延续了第一个函数芽，并到达原收敛圆盘之外。这里没有改变原级数的收敛半径；改变的是中心和所用局部表达式。
\end{example}

\section{沿路径的解析延拓及其唯一性}

设 $\gamma:[0,1]\to D$ 连续、$\gamma(0)=a$。沿 $\gamma$ 延拓函数芽，是给每个时刻 $t$ 指定一个位于 $\gamma(t)$ 的函数芽，并要求在每个时刻附近，这些芽都由同一个局部全纯函数表示。由区间紧性，可用有限分段描述：选
\[
 0=t_0<t_1<\cdots<t_N=1,
\]
使每段路径包含于一个开集 $U_j$，并在其上有全纯函数 $f_j$；相邻的 $f_j,f_{j+1}$ 在转换点 $\gamma(t_j)$ 的一个邻域内相同，起始芽为给定芽。必要时缩小邻域并细分区间，可使各段路径紧含于相应 $U_j$，转换点位于一个相容的交叠小圆盘中。

\begin{proposition}[固定路径上的唯一性]
\label{prop:ca-continuation-unique}
给定起始函数芽和一条路径，若沿该路径的解析延拓存在，则每个时刻的函数芽都唯一。
\end{proposition}
\begin{proof}
比较两次延拓，令 $E\subset[0,1]$ 为两芽相同的时刻集合。它包含零。任取 $t$，在其附近，两次延拓分别由 $\gamma(t)$ 周围的两个全纯函数表示，缩小到共同圆盘，并使一小段路径留在该圆盘内。若两芽在这段中的任一点相同，恒等定理便使两个函数在共同圆盘上相同，从而整小段上的芽相同。因此 $E$ 和其补集都相对开。区间连通，得到 $E=[0,1]$。
\end{proof}

固定路径上的唯一性不等于只由终点决定。例如从 $1$ 附近满足 $L(1)=0$ 的对数芽出发，沿 $\gamma(t)=e^{2\pi it}$ 延拓。每个局部分支都有导数 $1/z$，由链式法则，沿路径的分支值满足
\[
 \frac{d}{dt}L(\gamma(t))=2\pi i,
\]
所以回到 $1$ 后得到的值是 $2\pi i$，而非零。绕 $k$ 圈后，对数芽增加 $2k\pi i$，由其一半取指数得到的平方根芽则乘以 $(-1)^k$。

\begin{theorem}[单值化定理，monodromy theorem]
\label{thm:ca-monodromy}
设 $D$ 是区域，给定 $a\in D$ 处的全纯函数芽，并假设它可沿 $D$ 内从 $a$ 出发的每条连续路径解析延拓。若两条从 $a$ 到 $b$ 的路径在 $D$ 内定端点同伦，则终点函数芽相同。特别地，若 $D$ 单连通，该芽唯一确定 $D$ 上的单值全纯函数。
\end{theorem}
\begin{proof}
设同伦为 $H(t,s)$，对每个 $s\in[0,1]$，路径 $\gamma_s(t)=H(t,s)$ 从 $a$ 到 $b$。固定 $s_0$，由假设沿 $\gamma_{s_0}$ 存在延拓，选上一节中的有限相容函数链，使每段路径紧含于相应开集、每个转换点落在函数相等的交叠小圆盘内。紧集到这些开集补集的距离有正下界；$H$ 的一致连续性保证，当 $s$ 足够接近 $s_0$ 时，$\gamma_s$ 的每段仍落在同一开集，转换点仍落在同一交叠小圆盘。于是同一函数链也构成沿 $\gamma_s$ 的延拓。起点和终点固定，所以这些邻近路径的终点芽完全相同。

由命题~\ref{prop:ca-continuation-unique}，这就是各路径的唯一延拓。因此终点芽关于 $s$ 局部常值，在连通区间上必为常值，得到第一结论。若 $D$ 单连通，任意两条同端点路径定端点同伦：把第一条与第二条的反向拼成闭路，利用闭路收缩即可得到这种同伦。故可把 $F(b)$ 定义为任意路径延拓到 $b$ 后的值。对 $b$ 附近的小圆盘，先沿固定路径到 $b$，再沿圆盘内短线段移动，所得芽都由同一个局部函数表示，所以 $F$ 在该圆盘上全纯。各点均如此，$F$ 在 $D$ 全纯；唯一性由恒等定理得到。
\end{proof}

定理中的“可沿每条路径延拓”不能由单连通性替代。例如原点处的 $1/(1-z)$ 函数芽不能沿穿过 $1$ 的路径在整个复平面中全纯延拓，尽管复平面单连通。

\section{施瓦茨反射原理}

\begin{theorem}[施瓦茨反射原理，Schwarz reflection principle]
\label{thm:ca-reflection}
设 $B=B(x_0,r)$，其中 $x_0\in\mathbb R$。若 $f$ 在上半圆盘 $B^+=B\cap\{\operatorname{Im}z>0\}$ 全纯，并连续延拓到直径区间 $I=B\cap\mathbb R$，且 $f(x)\in\mathbb R$ 对 $x\in I$ 成立，则
\[
 F(z)=\begin{cases}
 f(z),&\operatorname{Im}z\ge0,\\
 \overline{f(\overline z)},&\operatorname{Im}z<0
 \end{cases}
\]
在 $B$ 上全纯。
\end{theorem}
\begin{proof}
下半圆盘内，由差商直接得导数 $F'(z)=\overline{f'(\overline z)}$。实轴处的实值边界条件使上下定义连续相接，因此 $F$ 在整个圆盘连续。对任意闭三角形 $T\subset B$，若它不跨实轴，沿与实轴平行的小平移先避开直径，再用柯西定理和连续性取极限，得到边界积分为零。若它跨实轴，按实轴将其切成上、下两个多边形；对各多边形采用同样的小平移论证，得到各自的边界积分为零。共有的实轴线段方向相反，积分抵消，故 $\int_{\partial T}F\,dz=0$。平移极限可以交换积分，是因为三角形附近的一个紧集上 $F$ 一致连续，边长又一致有界。莫雷拉定理遂给出全纯性。
\end{proof}

\begin{example}[边界条件决定反射方式]
在上半平面取 $f(z)=z^2+1$，实轴上取实值，反射得到同一个整函数。若取 $f(z)=iz$，实轴上的值通常为纯虚数，直接使用上一公式会在实轴产生不连续。此时应先对 $-if$ 使用实值反射，再乘回 $i$，得到下半平面公式 $F(z)=-\overline{f(\overline z)}=iz$。这说明反射公式必须由实际的边界值约束推导。
\end{example}

\section{自然边界：解析延拓可能处处受阻}

\begin{definition}[自然边界]
若 $f$ 在圆盘中全纯，却不能通过圆周的任何一点作局部全纯延拓，则该圆周称为 $f$ 的\emph{自然边界（natural boundary）}。通过边界点延拓，指存在该点的一个圆盘及其上的全纯函数，在与原圆盘相交的一侧和 $f$ 相同。
\end{definition}

\begin{example}[不用间隙定理的自然边界例子]
\label{ex:ca-natural-boundary}
函数
\[
 F(z)=\sum_{n=1}^{\infty}z^{n!}
\]
在单位圆盘内全纯，单位圆是它的自然边界。
\end{example}
\begin{solve}
当 $|z|\le r<1$ 时，$|z|^{n!}\le r^n$，故级数局部绝对一致收敛。设 $\zeta$ 为任意 $q$ 次单位根；当 $n\ge q$ 时，$\zeta^{n!}=1$。于是
\[
 F(r\zeta)=\sum_{n<q}r^{n!}\zeta^{n!}+\sum_{n\ge q}r^{n!}.
\]
第一项有界，第二项随 $r\uparrow1$ 趋于正无穷：给定任意正整数 $N$，先保留 $N$ 项，再令 $r$ 足够接近一，就能使这 $N$ 项的和大于 $N/2$。因此 $F$ 不能通过任何单位根全纯延拓。单位根在圆周上稠密，若能通过某个圆周点延拓，其延拓圆盘会包含附近一个单位根，并使该单位根附近的径向值有界，矛盾。故整个单位圆都是自然边界。
\end{solve}

这一例子也说明，圆周上的障碍可能是稠密的一整圈，而不是有限个孤立奇点；不能把所有幂级数的边界行为都归为极点或本性孤立奇点。

\section{黎曼曲面的定义与两个模型}

\begin{definition}[黎曼曲面]
\emph{黎曼曲面（Riemann surface）}是一个连通、豪斯多夫、第二可数的拓扑空间 $X$，每点有邻域 $U$ 及到复平面开集的同胚 $\varphi:U\to V\subset\mathbb C$，且任意两个这样的坐标在交叠部分的转移映射 $\varphi_2\circ\varphi_1^{-1}$ 双全纯。豪斯多夫指不同点可用不交开邻域分离；第二可数指拓扑有可数基。函数 $F:X\to\mathbb C$ 全纯，指每个局部表示 $F\circ\varphi^{-1}$ 全纯。
\end{definition}

坐标相容性保证全纯性的定义与所选坐标无关。复平面开区域和黎曼球面都是例子：球面在无穷远处用 $w=1/z$ 作坐标，交叠上的转移函数全纯且可逆。这个定义并不要求曲面作为三维空间中的光滑图形出现。

对数模型取
\[
 X_{\log}=\{(z,w)\in(\mathbb C\setminus\{0\})\times\mathbb C:e^w=z\}.
\]
映射 $w\mapsto(e^w,w)$ 把 $\mathbb C$ 同胚且双全纯地参数化为这个曲面，因此 $w$ 本身是全局坐标。投影 $\pi(z,w)=z$ 在每个非零 $z$ 上有 $w+2k\pi i$ 这些不同原像。局部选取对数使投影成为局部双全纯映射，而曲面上的 $L(z,w)=w$ 是单值全纯函数。它没有在去心平面上制造全局对数，因为同一个底面点仍有无穷多个提升点。

平方根模型取
\[
 X_{\sqrt{\phantom z}}=\{(z,w)\in\mathbb C^2:w^2=z\}.
\]
全局参数 $w\mapsto(w^2,w)$ 表明曲面本身仍与 $\mathbb C$ 双全纯等价。投影在 $w\ne0$ 时有两层局部逆，而在 $w=0$ 时导数 $d(w^2)/dw=0$，所以投影在此分支。曲面点 $(0,0)$ 有正常的局部坐标 $w$，并不是曲面本身的奇点。区分“曲面上的函数单值”“向底面的投影有分支”和“曲面自身奇异”，可以避免把不同对象混为一谈。

\section{习题与解答}

\begin{enumerate}
 \item 从 $1$ 处取值为 $1$ 的平方根芽出发，分别沿绕原点一圈、两圈的闭路延拓，求返回值。
 \item 若 $f$ 是圆盘上满足 $f(\overline z)=\overline{f(z)}$ 的全纯函数，证明其原点泰勒系数均为实数。圆盘以原点为中心。
 \item 对曲面 $X=\{(z,w):w^3=z\}$，指出一个全局坐标、投影的分支点，以及非零底面点的原像数。
\end{enumerate}
\begin{solve}[解答]
第一题返回值依次为 $-1$、$1$，因为每圈使连续辐角增加 $2\pi$，平方根的辐角增加 $\pi$。第二题把 $f(z)=\sum a_nz^n$ 代入对称关系，得到 $\sum a_n\overline z^{\,n}=\sum\overline{a_n}\overline z^{\,n}$；由幂级数唯一性，$a_n=\overline{a_n}$。第三题取坐标 $w$，投影为 $w\mapsto w^3$，仅在 $w=0$ 分支；每个非零底面点恰有三个原像。曲面本身由 $w\mapsto(w^3,w)$ 参数化，没有奇点。
\end{solve}

